% introduction.tex — Section 1: Introduction

\section{Introduction}
\label{sec:introduction}

\subsection{The Copernican program is not finished}
\label{subsec:copernican}

Physics has a recurring pattern.  Someone draws a line and declares
one side privileged.  The line holds for a while, then falls, and the
theory that remains is simpler and more predictive than the one it
replaced.

Copernicus removed the privileged position: Earth is not the center of
the universe~\citep{copernicus1543}.  Galileo removed the privileged
velocity: there is no absolute rest.  Einstein removed the privileged
inertial frame: the laws of physics take the same form for all
inertial observers~\citep{einstein1905}.  General relativity went
further: no coordinate system is privileged.  At each step, the
person who insisted ``but THIS distinction is real'' lost.

One distinction remains.  Contemporary physics draws a line between
mathematical structures that are ``physically realized'' and those
that are ``merely mathematical.''  The Standard Model, its coupling
constants, the dimensionality of spacetime --- these are declared
physical.  Other consistent mathematical structures --- an $\SU{5}$
gauge theory, a 7+1-dimensional spacetime, a universe governed by
quaternionic quantum mechanics --- are declared unphysical.  The line
is drawn, and one side is called reality.

This paper argues that removing the last distinction completes the
Copernican program.  \emph{No instantiation of a consistent
mathematical structure is ontologically privileged over any other.}
We call this principle \emph{radical relativity}.  It is not an axiom
being added to physics.  It is the refusal to add the axiom that some
structures are special.

The claim is not merely philosophical.  If radical relativity holds,
the question ``why is there something rather than nothing?'' is
malformed.  There is everything.  The productive question is: \emph{why
do observers experience this specific physics?}  We summarize a
derivation chain --- developed across seven companion
papers~\citep{ehrlich2026measure,ehrlich2026bb,ehrlich2026lipschitz,ehrlich2026born,ehrlich2026qm,ehrlich2026gr,ehrlich2026sm}
--- that answers this question.  The physics observers experience is
forced by the algebraic consequences of self-modeling.

\subsection{Summary of results}
\label{subsec:summary}

The derivation chain is the central artifact of this paper.
Table~\ref{tab:chain} presents it in full.  The chain begins with the
philosophical premise (radical relativity) and terminates at the
Standard Model gauge group with chiral fermion representation,
passing through complex quantum mechanics and general relativity.

\begin{table*}[t]
\centering
\caption{The derivation chain from radical relativity to the Standard
  Model.  ``Premise'' = philosophical starting point.  ``Framework'' =
  operational definition.  ``Theorem$^*$'' = mathematically proved,
  conditional on four structural assumptions (A1--A4, see text).
  ``Conditional'' = depends on identified gaps.}
\label{tab:chain}
\small
\begin{tabular}{@{}cp{4.5cm}p{3cm}p{3.5cm}@{}}
\toprule
Step & Claim & Status & Source \\
\midrule
0 & No privileged instantiation & Premise & This paper \\
1 & All consistent structures exist & Tautology under Step~0
  & \citet{tegmark2008} \\
2 & $\rho$ selects self-modelers & Framework
  & \citet{ehrlich2026measure} \\
3 & Self-modeling $\to$ Jordan algebras & \textbf{Theorem$^*$}
  & \citet{ehrlich2026qm} \\
4 & Faithful tracking $\to$ local tomography & \textbf{Theorem$^*$}
  & \citet{ehrlich2026qm} \\
5 & LT $\to$ complex $C^*$-algebras
  & \textbf{Theorem$^*$} & \citet{ehrlich2026qm} \\
6 & $C^*$-algebras $\to$ QM (Born, unitarity)
  & \textbf{Theorem}
  & \citet{gleason1957}, \citet{barandes2025} \\
7 & Locality $\to$ area law $\to$ GR
  & Conditional (4 gaps)
  & \citet{ehrlich2026gr} \\
8 & Non-composability $\to$ $\hthree$ $\to$ SM
  & Conditional (Gap~C)
  & \citet{ehrlich2026sm} \\
9 & $\rho$ resolves to unique $\rhoJ$ in $\hthree$
  & \textbf{Theorem}
  & \citet{ehrlich2026measure} \\
\bottomrule
\end{tabular}
\end{table*}

The chain's load-bearing joint is Steps~3--5.  These are proved from
one operational premise (faithful self-modeling) plus four structural
assumptions:
\begin{itemize}[nosep]
\item[(A1)] \emph{Finite-dimensionality.}  The state space has finite
  dimension.  Structural consequence: a self-modeler embedded in a
  finite region has finite information capacity (concretely, the
  Bekenstein bound gives an upper bound on the number of distinguishable
  states).
\item[(A2)] \emph{Faithfulness.}  The tracking map $\varphi$ is an
  order isomorphism.  This is the operational premise itself: the
  self-model preserves all operationally accessible structure.
\item[(A3)] \emph{Minimal composite.}  The body-model composite
  $V_{BM}$ is the smallest state space carrying product states,
  product effects, non-signaling constraints, and a product-form
  sequential product.  In plain language: the composite carries exactly
  the correlations that faithful tracking creates, and no more.  Extra
  composite structure beyond what product measurements require would be
  dead weight --- degrees of freedom that contribute nothing to
  tracking fidelity.  This is the assumption doing the most
  independent work: it is what forces $\dim(V_{BM}) = \dim(V_B) \cdot
  \dim(V_M)$ (local tomography), which in turn selects the complex
  field.  For complex quantum mechanics, the minimal and maximal
  composites coincide~\citep{bgw2020}, so the assumption is not
  restrictive for the theory we derive --- but it is load-bearing in
  the derivation.
\item[(A4)] \emph{Simple EJA.}  The algebra has no nontrivial
  ideal --- it cannot be decomposed into independent blocks.  A
  self-modeler that decomposes into independent subsystems is not
  modeling \emph{all} of itself: each block models only itself, not
  the whole.  Simplicity is the algebraic expression of ``whole models
  whole.''
\end{itemize}
No assumption about the complex field, Hilbert space structure, or
$C^*$-algebraic operations enters the premises; all emerge as
consequences.  This is the content of companion
Paper~5~\citep{ehrlich2026qm}.

Steps 7 and 8 are conditional.  The gravity derivation (Step~7)
inherits four standard gaps shared by all lattice quantum gravity
programs: emergent Lorentz invariance, lattice Bisognano-Wichmann,
local equilibrium, and the continuum limit.  The Standard Model
derivation (Step~8) is conditional on Gap~C, a conjectural
complexification principle whose resolution may require evolutionary
rather than algebraic arguments (Sec.~\ref{subsec:obj-gaps}).

\subsection{What this paper is and is not}
\label{subsec:scope}

This paper \textbf{is}:
\begin{itemize}[nosep]
  \item A philosophical manifesto for radical relativity.
  \item A roadmap to the companion papers (Papers 1--7) that contain
    the proofs and derivations.
  \item An honest register of what is proved, what is conditional,
    and what is open.
\end{itemize}

\noindent This paper \textbf{is not}:
\begin{itemize}[nosep]
  \item A source of new theorems.  All technical results appear in
    the companions.
  \item A theory of consciousness.  The framework makes a
    \emph{locative} claim: consciousness lives in self-referential
    fixed points.  Whether fixed points \emph{are} consciousness or
    merely \emph{host} it is not addressed.
  \item Theology.  The framework has implications for how religious
    traditions encode structural truths, but those implications are
    developed elsewhere.
\end{itemize}
